\documentclass[12pt]{article}

\usepackage{amsmath}
\usepackage{amssymb}
\usepackage{amsthm}
\usepackage[margin=1in]{geometry}
\usepackage{hyperref}
\usepackage{booktabs}
\usepackage{url}

\newtheorem{definition}{Definition}
\newtheorem{theorem}{Theorem}
\newtheorem{proposition}{Proposition}

\begin{document}

\begin{titlepage}
\centering
\vspace*{2cm}

{\LARGE\bfseries Fermi++: A Unified Theory of Representational Consistency, Quantum Normative Dynamics, and Pragmatic AI Verification}

\vspace{1.5cm}

{\Large Andrew H. Bond, Senior Member, IEEE}

\vspace{0.5cm}

{\large
Department of Computer Engineering\\
San Jos\'e State University\\
San Jos\'e, CA 95192, USA\\[0.3cm]
\texttt{andrew.bond@sjsu.edu}\\[0.3cm]
ORCID: 0009-0003-2599-6158
}

\vspace{1cm}

{\large Spring 2026}

\vspace{2cm}

\begin{abstract}
We present \textbf{Fermi++}, a unified formal framework for the verification of ethical and logical consistency in autonomous systems. Fermi++ integrates three previously disparate domains: (1) a categorical foundation for representational invariance using double categories, (2) a Noetherian field theory that derives moral conservation laws from symmetries of description, and (3) a quantum normative dynamics (QND) that models moral ambiguity via Hilbert space formalisms. We demonstrate that the ``Bond Index'' ($\beta$) serves as the universal coupling constant across these layers. Finally, we provide an instrumentalist justification for the framework, treating Fermi++ not as a metaphysical claim about the nature of morality, but as a robust engineering instrument for the ``Fermi-level'' verification of foundation models.
\end{abstract}

\vfill

\noindent\rule{\linewidth}{0.4pt}
\small
\noindent\textbf{AI-Assisted Writing Disclosure:} Portions of this manuscript were drafted and refined with AI assistance (Claude, Anthropic). All technical content, mathematical results, and claims were conceived, verified, and approved by the author.

\noindent\textbf{Ethics Statement:} This research involves no human subjects. Empirical validation uses publicly available corpora and synthetic data. No IRB approval was required.

\noindent\textbf{Data Availability:} Code and data are available at \url{https://github.com/ahb-sjsu/erisml-lib}.

\noindent\textbf{Conflict of Interest:} The author declares no conflicts of interest. This research received no external funding.

\noindent\textbf{CRediT Statement:} Andrew H. Bond: Conceptualization, Methodology, Formal Analysis, Investigation, Writing --- Original Draft, Writing --- Review \& Editing, Software.
\noindent\rule{\linewidth}{0.4pt}
\normalsize

\end{titlepage}

\noindent\textbf{Keywords:} AI safety, representational consistency, double categories, Noether's theorem, quantum normative dynamics, ethical field theory, foundation model verification

\bigskip
\hrule
\bigskip

\section{Introduction: The Fermi Gap}

In the era of large-scale foundation models, a ``Fermi Paradox'' of AI safety has emerged: despite the vast increase in computational intelligence, the ``signal'' of coherent moral reasoning remains elusive. Existing alignment techniques---such as RLHF---often produce ``surface consistency'' that collapses under re-description.

\textbf{Fermi++} is proposed as a multi-layered solution to this gap. It treats consistency not as a set of rules to be followed, but as a \textbf{structural invariant} of the system's internal representational space. The broader mathematical foundations for this approach are developed fully in~\cite{bond2026ge}.

\bigskip
\hrule
\bigskip

\section{The Categorical Foundation (The Engineering Layer)}

At its base, Fermi++ relies on the \textbf{Alignment Double Category} $\mathbb{A}$.

\begin{itemize}
    \item \textbf{Horizontal Morphisms ($\rightarrow$):} Represent \textit{re-descriptions} or fiber moves (e.g., paraphrasing an input).
    \item \textbf{Vertical Morphisms ($\downarrow$):} Represent \textit{scenario perturbations} or base moves (e.g., changing the actor's gender or the setting).
    \item \textbf{The Coherence Square:} For any input $x$, the diagram formed by $\rightarrow$ and $\downarrow$ must commute.
\end{itemize}

We define the \textbf{Bond Index ($\beta$)} as the integrated defect of these squares across the model's manifold. In Fermi++, $\beta$ is the ``temperature'' of the system's incoherence.

\bigskip
\hrule
\bigskip

\section{Ethicodynamics and Noether's Theorem (The Field Layer)}

Fermi++ posits that if a system is invariant under re-description (a symmetry), there must exist a conserved quantity. We identify this as \textbf{Harm Current ($J^\mu$)}.

Using the Lagrangian density $\mathcal{L}$, we derive the \textbf{Ethical Field Strength Tensor} $F_{\mu\nu}$:
\[
    F_{\mu\nu} = \partial_\mu A_\nu - \partial_\nu A_\mu
\]
where $A_\mu$ is the ``coherence potential.'' In this regime:

\begin{itemize}
    \item \textbf{Gauge Removable Defects:} Are local ``implementation glitches'' that can be smoothed via fine-tuning.
    \item \textbf{Intrinsic Anomalies:} Represent non-trivial curvature in the ethical manifold---logical contradictions in the system's core specifications.
\end{itemize}

\bigskip
\hrule
\bigskip

\section{Quantum Normative Dynamics (The Complexity Layer)}

As models move into domains of high moral ambiguity, classical ``Go/No-Go'' logic fails. Fermi++ incorporates \textbf{Quantum Normative Dynamics (QND)} to model these states:

\begin{itemize}
    \item \textbf{The Ethon ($\eta$):} The discrete carrier of moral influence between agents in a multi-agent system.
    \item \textbf{Superposition:} Ambiguous prompts are represented as a state $|\psi\rangle = \alpha|0\rangle + \beta|1\rangle$. The model's output is a ``measurement'' that triggers decoherence.
    \item \textbf{Entanglement:} Fermi++ accounts for non-local responsibility, where the moral status of Agent A cannot be evaluated without reference to its entanglement with Agent B.
\end{itemize}

\bigskip
\hrule
\bigskip

\section{The Pragmatist Grounding}

Crucially, Fermi++ rejects ``Moral Realism.'' Following the \textbf{Pragmatist Rebuttal} (Bond, 2025), we treat the Hilbert spaces and Tensors of Fermi++ as \textbf{adaptive instruments}.

The ``Ground'' of Fermi++ is not a divine or natural law, but the \textbf{Fermi Level of Coherence}: the highest state of complexity an AI can occupy before it begins to generate ``Harm Current'' through self-contradiction. We argue that ``Grace'' in an AI system is manifested as favorable ``vacuum fluctuations''---emergent beneficial behaviors that arise from a highly coherent internal geometry.

\bigskip
\hrule
\bigskip

\section{Implementation: The erisml-lib v0.3.0}

The theoretical stack of Fermi++ is operationalized via the erisml-lib, now upgraded to DEME 3.0 with tensorial ethics support~\cite{bond2026erisml}. By calculating Rank-4 Tensors:
\[
    T^{\alpha\beta\gamma\delta}
\]
the library computes the \textbf{Commutator Defect ($\Delta_C$)} and the \textbf{Mixed Defect ($\Delta_M$)}, providing a real-time ``Coherence Dashboard'' for model deployment.

\bigskip
\hrule
\bigskip

\section{Conclusion: Beyond the Paradox}

Fermi++ provides the first unified path from \textbf{abstract category theory} to \textbf{HPC-scale AI testing}. By treating ethics as a measurable physical field and a verifiable categorical structure, we move AI alignment from a philosophical debate to an engineering discipline. The Fermi Paradox of AI is solved not by finding ``aliens'' (AGI), but by ensuring that the intelligence we do build is structurally incapable of the incoherence that leads to systemic harm.

\bigskip
\hrule
\bigskip

\begin{thebibliography}{9}

\bibitem{bond2025rep}
Bond, A.~H. (2025).
\textit{Representational Consistency for Machine Learning Systems}.
IEEE TAI.

\bibitem{bond2025noether}
Bond, A.~H. (2025).
\textit{Noether's Theorem for Ethics}.

\bibitem{bond2025qnd}
Bond, A.~H. (2025).
\textit{Quantum Normative Dynamics}.

\bibitem{bond2025pragmatist}
Bond, A.~H. (2025).
\textit{A Pragmatist Rebuttal to Logical and Metaphysical Arguments}.

\bibitem{bond2026ge}
A.~H.~Bond, \emph{Geometric Ethics: The Mathematical Structure of Moral Reasoning}, v1.14, Department of Computer Engineering, San Jos\'e State University, Spring 2026.

\bibitem{bond2026erisml}
A.~H.~Bond, ``ErisML/DEME v3.0: A Library for Governed, Foundation-Model-Enabled Agents with Tensorial Ethics,'' \url{https://github.com/ahb-sjsu/erisml-lib}, 2026.

\end{thebibliography}

\end{document}
